\documentclass[11pt]{article}

\usepackage[margin=1in]{geometry}
\usepackage{array}
\usepackage{booktabs}
\usepackage{enumitem}
\usepackage{hyperref}
\usepackage{longtable}
\usepackage{tabularx}
\usepackage{xcolor}

\hypersetup{
  colorlinks=true,
  linkcolor=blue!55!black,
  urlcolor=blue!55!black,
  citecolor=blue!55!black
}

\emergencystretch=3em
\newcommand{\code}[1]{\texttt{\detokenize{#1}}}
\newcommand{\grip}{\code{grip}}
\newcommand{\nd}{ND}

\title{Weighted GRIP ND Layout Implementation}
\author{grip design note}
\date{2026-05-19}

\begin{document}

\maketitle

\section{Purpose and Scope}

This note describes the dimension-general weighted GRIP layout implementation in
the \grip{} R package. The implementation goal is deliberately conservative:
\code{grip.layout.weighted()} should use the dimension-general backend while
reconstructing the historical weighted GRIP algorithm exposed by
\code{grip.layout.weighted.legacy()} and
\code{grip.layout.globalrep.weighted()}. The \nd{} backend replaces
hard-coded two- and three-dimensional coordinate logic with loops over
\code{dim >= 2}; it is not intended to define a new weighted layout algorithm.

The document is implementation-facing. It records the public call surface, the
legacy-to-\nd{} mapping, the multilevel algorithm, the parity-sensitive
arithmetic and state updates, the diagnostic paths, and the validation evidence
available after the weighted-\nd{} parity campaign.

\section{Public API}

The current weighted public API has two roles:

\begin{itemize}[leftmargin=*]
  \item \code{grip.layout.weighted()} is the default weighted layout wrapper and
  forwards to \code{grip.layout.weighted.nd()}.
  \item \code{grip.layout.weighted.legacy()} preserves the historical
  compiled weighted backend and forwards to
  \code{grip.layout.globalrep.weighted()}.
\end{itemize}

The two public wrappers are signature-aligned. The parity checkpoint on
2026-05-19 compared \code{formals(grip.layout.weighted.legacy())} and
\code{formals(grip.layout.weighted())} and found no legacy-only arguments, no
\nd{}-only arguments, and no default-value differences.

\subsection{Main Public Weighted Parameters}

The weighted wrappers accept graph input either as an edge matrix with optional
\code{edge_weights}, or as parallel \code{adj_list} and \code{weight_list}
objects. The weighted edge values are interpreted as desired graph-metric edge
lengths and are normalized by \code{length_normalization}, whose choices are
\code{"median"}, \code{"mean"}, and \code{"none"}.

The main control groups are:

\begin{itemize}[leftmargin=*]
  \item multilevel schedule: \code{rounds}, \code{final_rounds},
  \code{num_init}, \code{num_nbrs};
  \item local temperature: \code{r}, \code{s}, \code{tinit_factor};
  \item weighted force scaling: \code{repulsion_factor},
  \code{coarse_repulsion_factor}, \code{coarse_repulsion_sample},
  \code{coarse_repulsion_exact_below};
  \item final-stage behavior: \code{final_anchor_factor},
  \code{final_move_scale_after_first}, \code{final_mode};
  \item metric-neighbor search: \code{metric_neighbor_cap};
  \item insertion behavior: \code{insertion_anchor_count},
  \code{insertion_anchor_scope}, \code{insertion_anchor_strategy},
  \code{level0_insertion_mode}, \code{level0_anchor_count},
  \code{level0_local_kk_steps};
  \item optional LGKK refinement: \code{lgkk_polish_rounds},
  \code{lgkk_multiscale_rounds}, \code{lgkk_rounds_coarse},
  \code{lgkk_rounds_pre_final}, \code{lgkk_rounds_final},
  \code{lgkk_local_nbrs}, \code{lgkk_landmark_count},
  \code{lgkk_multiscale_scope}, and \code{lgkk_active_limit}.
\end{itemize}

\section{Call Paths}

\subsection{Final Layout Paths}

\begin{longtable}{p{0.31\textwidth}p{0.61\textwidth}}
\toprule
Path & Implementation \\
\midrule
\endhead
\code{grip.layout.weighted()} &
Default public weighted wrapper. Forwards to
\code{grip.layout.weighted.nd()} in \code{R/grip_layout_weighted_nd.R}. \\
\code{grip.layout.weighted.nd()} &
Validates weighted \nd{} inputs, resolves presets, normalizes lengths, handles
disconnected components, calls \code{grip_layout_weighted_nd_adj_cpp()}, and
optionally applies post-layout LGKK polish. \\
\code{grip_layout_weighted_nd_adj_cpp()} &
Rcpp entry point in \code{src/grip_nd_rcpp.cpp}. Converts R adjacency and
weights to \code{GraphND}, parses mode strings, constructs
\code{DrawGraphND}, and returns the coordinate matrix from
\code{DrawGraphND::layout()}. \\
\code{grip.layout.weighted.legacy()} &
Legacy public wrapper. Forwards to \code{grip.layout.globalrep.weighted()} in
\code{R/grip_layout_weighted.R}. \\
\code{grip.layout.globalrep.weighted()} &
Legacy weighted R path. Validates the same public controls, calls
\code{grip_layout_globalrep_weighted_adj_cpp()}, and optionally applies
post-layout LGKK polish. \\
\code{grip_layout_globalrep_weighted_adj_cpp()} &
Legacy Rcpp entry point in \code{src/grip_rcpp.cpp}. Constructs
\code{DrawGraph} with \code{weightedCore = true}, configures metric search, and
runs the legacy weighted engine. \\
\bottomrule
\end{longtable}

\subsection{Diagnostic Trace Paths}

The strict diagnostic comparisons use trace entry points rather than only final
layout wrappers:

\begin{itemize}[leftmargin=*]
  \item Legacy diagnostic path:
  \code{grip.layout.trace.weighted()} $\rightarrow$
  \code{grip_layout_globalrep_weighted_trace_adj_cpp()}.
  \item \nd{} diagnostic path:
  \code{grip.layout.weighted.nd.trace()} $\rightarrow$
  \code{grip_layout_weighted_nd_trace_adj_cpp()}.
\end{itemize}

Both trace paths record layout frames and metadata. The \nd{} trace path also
supports gated refinement-step diagnostics used during parity debugging. These
diagnostics are disabled by default and are test-only instrumentation, not
ordinary user behavior.

\section{Algorithm Overview}

Weighted GRIP is a multilevel layout algorithm driven by a maximal independent
set filtration (MISF). The weighted variant replaces unweighted graph distances
with weighted shortest-path distances derived from the edge weights.

The high-level order is:

\begin{enumerate}[leftmargin=*]
  \item validate R inputs and normalize weighted lengths;
  \item construct the weighted graph representation;
  \item build the weighted MISF and its active-vertex order;
  \item initialize the coarsest active set;
  \item descend levels from coarse to fine;
  \item insert newly active vertices using legacy-compatible anchors and
  initial-placement rules;
  \item refine the active set with weighted KK-style rounds on coarse levels;
  \item refine the full graph with weighted FR-style rounds, or with the
  alternate weighted KK-repulse final mode;
  \item optionally apply final anchor and final move-scaling controls;
  \item optionally run compiled multiscale LGKK rounds inside the multilevel
  solver and R-level post-layout LGKK polish after the main solve.
\end{enumerate}

\section{Legacy-to-ND Mapping}

\begin{longtable}{p{0.32\textwidth}p{0.32\textwidth}p{0.26\textwidth}}
\toprule
Legacy implementation & \nd{} implementation & Role \\
\midrule
\endhead
\code{R/grip_layout_weighted.R} &
\code{R/grip_layout_weighted_nd.R} &
Public weighted wrappers, validation, normalization, disconnected-component
handling, and R-level LGKK polish. \\
\code{grip_layout_globalrep_weighted_adj_cpp()} &
\code{grip_layout_weighted_nd_adj_cpp()} &
Compiled final-layout entry point. \\
\code{grip_layout_globalrep_weighted_trace_adj_cpp()} &
\code{grip_layout_weighted_nd_trace_adj_cpp()} &
Compiled trace entry point. \\
\code{DrawGraph} with \code{weightedCore = true} &
\code{DrawGraphND} &
Stateful multilevel layout engine. \\
\code{create_misf_weighted()} &
\code{build_weighted_misf_nd()} &
Weighted MISF construction and active ordering. \\
\code{compute_weighted_shortest_paths()} and
\code{metricNbrs[root][level]} &
\code{populate_metric_neighbors_for_root()} and
\code{metric_neighbors_cache_[root][level]} &
Cached weighted metric-neighbor layers. \\
\code{KK_spring_weighted_local_v1()} &
\code{local insertion polish in DrawGraphND} &
Local weighted insertion refinement. \\
\code{KK_spring_weighted_v1()} &
\code{legacy_weighted_kk_displacement()} &
Coarse-level weighted KK displacement. \\
\code{KK_spring_weighted_final_v1()} &
\code{legacy_weighted_kk_final_displacement()} &
Alternate full-graph KK-repulse final mode. \\
\code{FR_spring_weighted_v1()} &
\code{legacy_weighted_fr_displacement()} &
Default full-graph weighted FR displacement. \\
\code{add_final_anchor_force()} &
\code{add_final_anchor_force()} &
Optional final anchor term. \\
\code{mish_engine_weighted()} &
\code{DrawGraphND::layout()} and
\code{refine_legacy_weighted_level()} &
Top-level weighted multiscale schedule and per-round update loop. \\
\bottomrule
\end{longtable}

\section{Weighted MISF and Metric Neighbors}

The \nd{} MISF builder is \code{build_weighted_misf_nd()} in
\code{src/WeightedMisfND.cpp}. It mirrors the legacy weighted MISF structure:
the active order is stored in \code{WeightedMisfND::order}, vertex positions in
that order are stored in \code{inverse}, per-vertex depth is stored in
\code{vertex_depth}, and per-level active sizes are stored in
\code{level_size}.

Weighted shortest-path traversal uses a Dijkstra priority queue ordered first
by distance and then by vertex id. This tie-breaking is parity-sensitive
because MISF membership, metric-neighbor order, and insertion anchors all
depend on traversal order. The weighted radius schedule is
\code{2^(level - 1)} for levels greater than one and \code{1} for levels zero
and one.

Metric neighbors are cached by root and level. In the legacy code the cache is
\code{metricNbrs[root][level]}; in the \nd{} code it is
\code{metric_neighbors_cache_[root][level]}. The \nd{} path populates all
needed layers for a root on first use through
\code{populate_metric_neighbors_for_root()} and then retrieves the requested
layer through \code{metric_neighbors_for_level()}. This is intentionally a
cache-equivalent port rather than a fresh shortest-path query at every force
evaluation.

\section{Insertion and Initial Placement}

At each MISF level, newly introduced vertices are inserted before the active
set is refined. Anchor selection is controlled by
\code{insertion_anchor_count}, \code{insertion_anchor_scope}, and
\code{insertion_anchor_strategy}. The current \nd{} backend implements the
legacy strategies, including selection from any higher level or only the
previous MISF level.

Initial coordinates are controlled by \code{level0_insertion_mode} and the
legacy weighted placement helpers:

\begin{itemize}[leftmargin=*]
  \item \code{"inherit"} preserves the level's default legacy behavior;
  \item \code{"barycenter"} uses anchor barycenters;
  \item \code{"least_squares"} uses the weighted local least-squares initial
  placement path.
\end{itemize}

Disconnected or no-anchor insertion is also parity-sensitive. The legacy code
falls back to \code{rand_Point()}; the \nd{} backend preserves this behavior
through legacy-compatible random coordinate generation rather than placing the
vertex at the origin.

\section{Refinement Loop}

The parity-critical update loop is \code{DrawGraphND::refine_legacy_weighted_level()}.
For each round it iterates over the first \code{active_count} vertices in
\code{misf.order}. For each active vertex it:

\begin{enumerate}[leftmargin=*]
  \item reads the cached metric-neighbor layer for the current MISF level;
  \item computes the displacement using the phase-specific force kernel;
  \item records diagnostic pre-temperature state when requested;
  \item updates the local temperature;
  \item stores \code{old_disp} and \code{old_disp_norm};
  \item scales displacement by heat and by the pre-temperature displacement
  norm;
  \item applies \code{final_move_scale_after_first} in the default final FR
  mode after the first final round when the scale is less than one;
  \item after all vertices in the round have computed their displacements,
  applies coordinate updates in active order.
\end{enumerate}

The delayed coordinate application matters: all force computations in a round
see coordinates from the start of that round, matching the legacy engine's
round-level update semantics.

\section{Force Kernels}

\subsection{Coarse Weighted KK}

The coarse weighted KK displacement is implemented by
\code{legacy_weighted_kk_displacement()}. For each metric neighbor it computes
the vector from the active vertex to the neighbor, the squared Euclidean norm,
the desired Euclidean distance \code{neighbor.dist * kLegacyEdgeND}, and the
legacy scale \code{norm2 / desired2 - 1}. The scaled vector is accumulated into
\code{disp_[vert]} in dimension order. On coarse levels, active-set repulsion
is added after attraction and before displacement normalization.

\subsection{Default Full-Graph Weighted FR}

The default final phase is \code{legacy_weighted_fr_displacement()}. Its
attraction term iterates over graph adjacency neighbors, not metric-neighbor
layers. For an edge from \code{vert} to \code{overt} with weight \code{w}, the
desired length is \code{kLegacyEdgeND * w}; the attraction step is the edge
vector multiplied by \code{norm2 / desired2}. The metric-neighbor layer then
contributes the final-level repulsion term scaled by
\code{repulsion_factor * 0.05 * kLegacyEdgeND^2 / norm2}.

During parity work, this kernel was one of the most sensitive areas because
the legacy implementation used \code{Point<>} arithmetic and historical
floating-point casts. The \nd{} path was adjusted so the arithmetic sequence,
neighbor order, and displacement scaling match the legacy behavior for the
strict parity fixtures.

\subsection{Alternate Final KK-Repulse Mode}

When \code{final_mode = "kk_repulse"}, the \nd{} engine calls
\code{legacy_weighted_kk_final_displacement()} at MISF level zero. This keeps
the same weighted metric-neighbor KK attraction as the coarse phase but uses
the finest-level repulsion scale. This path is included for signature and
algorithmic parity with the legacy backend.

\subsection{Active-Set Repulsion}

\code{add_legacy_active_repulsion()} implements the public coarse-repulsion
controls. If the active set is at or below
\code{coarse_repulsion_exact_below}, or if the sample size covers the active
population, it uses exact active-set repulsion. Otherwise it samples active
vertices with the legacy RNG and rescales by
\code{(active_count - 1) / sample_count}. The exact and sampled branches both
accumulate in active-order dimension loops.

\subsection{Final Anchor Force}

The final anchor path has two pieces. \code{prepare_final_anchors()} stores
coordinates for the active vertices before the final phase when
\code{final_anchor_factor > 0}. \code{add_final_anchor_force()} then adds
\[
  (\mathrm{anchor}_v - x_v) \cdot
  \frac{\code{final_anchor_factor}}{\code{kLegacyEdgeND}}
\]
to each dimension of the displacement. When
\code{final_anchor_factor = 0}, the function returns before mutating
displacement state.

\section{Dimensional Generalization}

The \nd{} backend stores coordinates as \code{PointND} values and loops over
\code{d = 0, \ldots, dim - 1}. The intended generalization is mechanical:

\begin{itemize}[leftmargin=*]
  \item replace explicit \code{x}, \code{y}, and \code{z} fields with indexed
  coordinate access;
  \item preserve neighbor and active-vertex ordering;
  \item preserve force accumulation order across neighbors and dimensions;
  \item preserve the local temperature sequence and delayed coordinate update;
  \item extend valid dimensions from legacy \code{dim \in \{2, 3\}} to
  \code{dim >= 2}.
\end{itemize}

This generalization is intentional and should not be considered an algorithmic
difference by itself. Algorithmic differences are only those that alter the
legacy schedule, force equations, state updates, random choices, cache
behavior, or public option semantics.

\section{LGKK Refinement}

Weighted \nd{} supports both R-level post-layout LGKK polish and compiled
multiscale LGKK stages inside the multilevel solver. The compiled stage is
controlled by per-stage budgets:
\code{lgkk_rounds_coarse}, \code{lgkk_rounds_pre_final}, and
\code{lgkk_rounds_final}, with \code{lgkk_multiscale_rounds} as the legacy
shared fallback. The scope is controlled by \code{lgkk_multiscale_scope}, and
\code{lgkk_active_limit} limits cache construction on large active sets.

The implementation builds sparse local and landmark graph-metric pairs on the
active set, applies LGKK displacement rounds, and uses the same local
temperature machinery as the main weighted refinement. The parity suite
includes compiled LGKK checks and staged-budget checks.

\section{Disconnected Components}

The public R wrappers detect disconnected graphs before calling the compiled
core. With \code{disconnected = "components"}, each component is induced,
laid out separately, and packed into one coordinate matrix. With
\code{disconnected = "error"}, the wrapper stops. The compiled weighted core
therefore continues to assume a connected graph, while the public wrappers
provide safe disconnected-graph behavior.

\section{Trace and Diagnostic Instrumentation}

The trace interfaces record frame matrices and metadata describing phase,
level index, MISF level, round number, and active vertex count. The
\nd{} trace backend also exposes optional refinement-step diagnostics used to
compare per-vertex and per-term force accumulation. These diagnostics were
central during parity repair because they could isolate differences in:

\begin{itemize}[leftmargin=*]
  \item active order and focal vertex order index;
  \item attraction edge identity, edge weight, desired length, and term vector;
  \item metric-neighbor identity, distance, and repulsion term;
  \item pre-temperature displacement and local heat update;
  \item post-temperature move applied to coordinates.
\end{itemize}

The refinement-step diagnostic is gated in tests. In the latest gflow run it
skipped because no divergent trace frame was found at tolerance \code{1e-08};
there was no focal divergence to expand.

\section{Validation Evidence}

The checkpoint note
\code{dev/design/weighted-grip/weighted_grip_nd_parity_checkpoint_2026-05-19.md}
records the current parity evidence. The key strict gates completed after the
LGKK port and formatting audit were:

\begin{longtable}{p{0.58\textwidth}p{0.32\textwidth}}
\toprule
Command or gate & Result \\
\midrule
\endhead
\code{devtools::test(filter = "layout-weighted-nd\$")} &
91 pass. \\
\code{GRIP_RUN_GFLOW_LGKK_PARITY_TESTS=true} with enforcement &
12 pass, 7 gated skips. \\
\code{devtools::test(filter = "layout-weighted")} &
153 pass, 10 gated skips. \\
\code{GRIP_RUN_GFLOW_FINAL_ANCHOR_PARITY_TESTS=true} with enforcement &
6 pass. \\
\code{GRIP_RUN_TRACE_PARITY_TESTS=true} &
13 pass, 1 gated skip. \\
\code{GRIP_RUN_FINAL_ANCHOR_TRACE_PARITY_TESTS=true} &
6 pass, 1 gated skip. \\
\code{GRIP_RUN_GFLOW_FULL_PARITY_TESTS=true} with enforcement &
13 pass, 7 gated skips. \\
\code{GRIP_RUN_GFLOW_STRESS_PARITY_TESTS=true} with enforcement &
13 pass, 7 gated skips. \\
\code{GRIP_RUN_GFLOW_TRACE_ALL_TESTS=true} with enforcement &
10 pass, 7 gated skips. \\
\code{GRIP_RUN_GFLOW_REFINEMENT_STEP_TRACE_TESTS=true} &
Skipped because no divergent gflow trace frame was found at tolerance
\code{1e-08}. \\
\code{make check-fast} &
Completed with known 2 warnings and 1 note. \\
\code{make check} &
Completed with known 2 warnings and 3 notes. \\
\bottomrule
\end{longtable}

The defensible claim is therefore tested full legacy-vs-\nd{} parity across
the current weighted parity suite, including the previously missing public
paths: coarse-repulsion controls, metric-neighbor cap, final anchor, final
move scaling, alternate final-stage mode, disconnected insertion behavior, and
LGKK stages.

\section{Known Limits of the Claim}

The parity claim is empirical, not a formal proof. It means that the current
strict tests, trace gates, and gflow-backed fixtures pass under enforcement.
It does not prove bitwise identity for every possible graph, seed, tuning
combination, compiler, or platform. In particular:

\begin{itemize}[leftmargin=*]
  \item floating-point behavior may still vary across platforms or compilers;
  \item parity is only as broad as the maintained fixture suite;
  \item legacy behavior is intentionally preserved even where it is historical
  rather than mathematically ideal;
  \item the package still has non-parity check warnings and notes, including
  non-portable compiler-warning flags, undeclared \code{geometry} usage in
  tests, CRAN new-submission status, a local time-verification note, and an
  HTML Tidy note.
\end{itemize}

\section{Maintenance Guidance}

Future changes to weighted layout code should preserve the distinction between
algorithmic modernization and parity preservation. A safe workflow is:

\begin{enumerate}[leftmargin=*]
  \item change legacy and \nd{} arithmetic together when the goal is
  modernization, so strict legacy-vs-\nd{} parity can remain enforced;
  \item record old-vs-new behavioral drift with layout snapshots when changing
  numerical semantics;
  \item rerun targeted parity gates for any touched phase before running broad
  gflow gates;
  \item update this design note when changing public options, backend force
  kernels, MISF construction, metric-neighbor cache semantics, insertion
  behavior, final-stage behavior, or LGKK staging.
\end{enumerate}

\section{Source Index}

\begin{itemize}[leftmargin=*]
  \item \code{R/grip_layout_weighted_nd.R}: \nd{} public weighted wrapper,
  default \code{grip.layout.weighted()} forwarding, trace wrapper, validation,
  disconnected handling, R-level LGKK polish.
  \item \code{R/grip_layout_weighted.R}: legacy weighted public wrapper and
  trace wrapper.
  \item \code{src/grip_nd_rcpp.cpp}: \nd{} Rcpp entry points and R-to-C++
  validation/parsing.
  \item \code{src/DrawGraphND.h} and \code{src/DrawGraphND.cpp}: \nd{}
  weighted layout engine, force kernels, insertion, tracing, final anchor,
  coarse repulsion, and LGKK stages.
  \item \code{src/WeightedMisfND.h} and \code{src/WeightedMisfND.cpp}:
  weighted MISF construction.
  \item \code{src/grip_rcpp.cpp}, \code{src/DrawGraph.h},
  \code{src/DrawGraph.cpp}, and \code{src/MishWeighted.cpp}: legacy weighted
  backend and reference behavior.
  \item \code{tests/testthat/helper-weighted-nd-trace-parity.R} and
  \code{tests/testthat/test-layout-weighted-nd-trace-parity.R}: trace parity
  helpers and tests.
  \item \code{tests/testthat/helper-weighted-nd-gflow-parity.R} and
  \code{tests/testthat/test-layout-weighted-nd-gflow-parity.R}: gflow-backed
  final-layout, stress, trace, LGKK, and refinement-step parity gates.
\end{itemize}

\end{document}
